\section{讨论}
\subsection{较明显的报告回升点估计位于高海拔高绿度环境}
报告回升的空间含义来自环境层内的变化，而非全岛平均值本身。低海拔低、高绿度组合从两次飓风之间到 Fiona 后分别增加 0.39 和 1.07 个百分点，高海拔高绿度组合增加 6.97 个百分点。后一组合同时具有较高的初期报告水平，因此其较大绝对增量应结合基线理解。低地两类组合末期仍低于初期的点估计，提示回升没有表现为所有环境同步返回此前水平。各层末期相对初期差值的空间聚类区间均跨越零，当前证据更明确地揭示了环境间的报告水平差异，尚不足以将环境层确定划分为已经恢复和未恢复。

这一解释还取决于哪些环境进入了分析。高海拔 NDVI 缺失率为 49.6\%，高海拔清单在缺失与完整子集中分别占 11.2\% 和 3.4\%。缺失记录相对完整记录的报告优势比在控制海拔与时期后由 1.86 降至 1.03，说明总体报告差异有明显的环境组成来源。标准化比较通过保持层内协变量分布一致，提高了时期对照的可比性；补充缺失地点的植被信息则能够扩大环境解释所覆盖的观测范围。两者共同服务于判断回升发生在哪里，避免将环境采样权重变化误读为所有地点共同变化。这一问题与 Backstrom 等\cite{backstrom2025}发现的采样偏差相互作用一致，说明控制调查努力之外，还需检查各环境层是否得到可比的代表。

海拔与绿度关联与波多黎各飓风后森林损伤及鸟类响应的生态背景相符\cite{uriarte2019,wunderle1995}。Howe 等\cite{howe2025}结合 LiDAR 与 NDVI 描述了红树林结构和绿度的互补变化，提示仅凭 NDVI 不能判定冠层结构或花资源已经恢复。Rivera-Milán 等\cite{rivera2025}发现波多黎各平鸽在飓风后仍长期处于低种群水平，也说明经历同一扰动的鸟类不一定具有相同的恢复轨迹。这些研究为本文提供生态背景，并非对绿蜂鸟结果的独立验证。海拔还同时对应气候、生境和可达性，三个时期包含两次风暴及其间的时间积累，缺少未受扰动的反事实，因而不能把期际差异单独归因于 Fiona。

\subsection{山地高报告与较高占域持续概率一致，低地内部存在绿度差异}
高海拔高绿度环境同时具有较高报告水平和较高的占域持续概率，两个结果在环境方向上相一致。校正观测排列后，海拔和绿度与较低局地灭绝概率的关联保留，而初期探测系数的绝对值缩小约 90\%。这保留了将较高、较绿地点的高报告水平与占域持续性联系起来的依据，同时削弱了原先较强的时期探测差异解释。报告模型与占域模型共享观测，这种一致性属于不同响应定义之间的相互参照，尚非独立的生物学验证。

低地结果使这一解释更具体。低海拔高绿度组合的末期报告预测为 5.17\%，低于初期的 6.27\%，但对应地点的平均占域持续概率为 71.3\%，高于低海拔低绿度地点的 24.2\%。因此，报告未完全回升与较高的占域持续概率能够并存。前者描述清单层面的报告轨迹，后者描述前期已占据地点的保留；两者的差别提示低地监测既要复访此前记录到该物种的地点，也要复访此前未记录到该物种的地点，并在不完全探测条件下估计状态转移。将所有低地地点视为同一种恢复状态，会掩盖绿度相关的内部差异。

固定海拔 100~m 时，NDVI 从 0.30 增至 0.76 对应占域持续概率增加 42.2 个百分点，而定殖变化仅为 0.7 个百分点且区间跨零。现有模型对绿度与既有地点保留的联系提供了较明确支持，对绿度促进新地点进入占据状态的支持较弱。等访问次数下，低地低、高绿度的再次报告分别为 1/9 和 15/28，也提供了相同方向的观测线索；这些小样本报告历史仍混合探测与状态变化。高绿度低地因而适合作为占域监测的重点复访层，低绿度低地则提供必要对照，以区分较低占域概率与探测不足对稀少报告的贡献。

这些过程概率描述当前模型尺度上的地点间关联。跨时期平均 NDVI 不能确定同一地点植被恢复的时间顺序，多年主要时期中的重复访问也依赖封闭假设。由此，当前结果帮助确定在哪些环境中设置重复调查网格，以估计占域持续概率；判断植被变化是否提高了持续性，需要在更短的合理调查窗口中同步记录地点状态与植被变化。Gordon 等\cite{gordon2024}利用重复定点调查，将定殖和灭绝与生境管理扰动后的时间联系起来。对于绿蜂鸟，在固定网格同步积累扰动与植被历史，有助于分别检验生境变化与两类状态转移的关系。

\subsection{监测应兼顾山地占域持续性与低地内部差异}
低地仅有 3.3\% 的清单报告该物种，却承载完整样本中 60.7\% 的正报告。它既是记录识别困难的环境，也是不能从持续性监测中省略的环境。随机森林在低、中海拔改善 AP，在高海拔没有同等改善，表明预测增益主要帮助从大量未报告清单中找到低地线索。全局前十分位中，低地正报告覆盖由 m4 的 26.4\% 增至 34.6\%，而山地覆盖下降。预测更准确并不自动带来更均衡的环境监测，模型排序需要服从环境对照的调查目标。

等预算比较给出了这一选择的可量化代价。全局随机森林排序选中 57.9\% 的高海拔清单和 6.9\% 的低海拔清单；各层约选 10\% 后，低地报告覆盖升至 45.6\%，总体覆盖却减少 388 条正报告。若目的是增加总体命中，全局排序更有利；若目的是比较低地不同条件下的占域持续概率，则需要主动保留低地名额。该实验比较的是历史清单筛查预算，尚未测量独立地点数、实际复访成本或持续性估计精度的提高。

据此，监测可由三个互补部分组成。第一，保留山地固定重复点，积累高报告环境中的探测历史，估计占域状态转移，并补充山地低绿度等样本较少的条件。第二，在低地同时保留低、中、高绿度地点，比较不同绿度条件下的报告变化和占域持续概率。第三，在各环境层内使用经过空间验证的排序寻找补充候选地点，同时保留固定点及未报告地点，使候选筛查不会替代完整的状态转移调查。所有重复访问应记录时长、距离、同期植被及可获得的花资源信息，以区分环境变化和探测机会。Ovaskainen 等\cite{ovaskainen2026}的鸟类监测系统采用固定调查路线与间隔记录，为结合预测建模和稳定观测条件提供了近期实例。本文据此强调，即使模型排序发生变化，也应维持可比较的固定重复调查。

\subsection{空间与时间验证决定复访线索的适用范围}
预测实验支持上述监测分工，但其作用强度随使用场景改变。随机森林相对 m4 的 AP 优势从随机划分的 0.143 缩小到空间划分的 0.033，说明熟悉地点组合中的高分不足以代表新地理区块中的收益。扩大区块后，部分灵活模型的优势进一步减弱。空间验证因而帮助判断候选地点排序在多大范围内仍有信息价值，并支持在扩大调查区域时保持环境层的固定覆盖\cite{roberts2017,valavi2019,ploton2020}。当前无距离缓冲的岛内验证仍不能确定所有未采样环境的适用范围。Koldasbayeva 和 Zaytsev\cite{spacetime2025}同样强调物种分布预测中的分块评价和时间外测试，其结果支持使验证情境与目标复访范围相匹配，但不能据此认定本文区块尺度已经消除全部空间依赖。

Fiona 后时间留出中，随机森林前十分位召回率由 m4 的 53.28\% 提高至 59.80\%，说明早期学习的关系仍能帮助识别后期报告。它为利用既有记录寻找后续复访线索提供支持，但尚未分别证明各低地绿度层的时间推广收益。空间分层回答哪些环境受益，时间留出回答总体报告识别能否延续到后期；联合解释二者有助于制定复访计划，不能替代同地点的重复探测记录与占域转移分析。Lovell 等\cite{lovell2025}在蝴蝶研究中发现，环境对占域与丰度的空间和时间作用在分布范围的不同部分并不一致。这为本文提供方法上的参照，应进一步在环境层内检验时间推广能力，而非假定总体留出收益均匀适用于全部环境。

其余实验进一步明确预测工具如何服务环境问题。EVI2 结果表明报告排序收益不限于 NDVI；校准结果区分候选排序与预期报告量估计；努力响应和变量重要性强调保持调查条件可比。m3 与 m4 接近的空间表现说明，描述时期--海拔关系的交互项未必提高新地点排序，生态解释与预测用途需各有证据\cite{elith2009}。神经流程的训练来源重叠削弱其高分支持新增地点监测的能力，其响应曲线和消融可用于检查拟合信息，却不能加强占域持续性的过程结论。地图与曲线最终应落实为可复访的环境对照，并结合重复调查和探测模型估计原有占域网格在后续时期仍被占据的概率。

\section{结论}
波多黎各绿蜂鸟的报告回升随海拔和植被绿度而异。高海拔高绿度环境的报告水平较高，绝对回升幅度的点估计也较大，但各环境层相对初期的恢复程度仍需重复调查确认。这些环境同时对应模型估计的较高占域持续概率；低地较绿地点也具有较高的占域持续概率，说明低地报告回升不足不能直接等同于所有原有占域网格都已不再被占据。预测改善主要帮助识别低、中海拔稀少报告，空间与时间验证则界定这些线索用于复访的范围。监测应保持山地固定重复点、低地绿度对照和层内候选地点筛查，分别分析报告变化、再次记录，以及考虑不完全探测后的占域持续概率和定殖概率。

\section*{待作者补充的声明}
\textbf{数据与代码可用性。}原始观测与协变量来自文献\cite{ortiz2026}的公开包。本修订包含源文件、保留的实验输出、证据审计及可重建汇总脚本。新增分析的公共归档地址与适用再分发条件仍需确认。\\
\textbf{作者、经费与利益冲突。}作者姓名、单位、贡献、经费及利益冲突声明需由作者提供。\\
\textbf{伦理。}本重分析未新增动物操作或野外干预。适用的机构认定与许可尚待确认，本文不声称已获得批准或豁免。\\
\FloatBarrier
